\section{Three heights and the explicit zero-free region}
\label{sec:three-heights}

We prove that \(c=2^{-20}\) is admissible. Suppose, to the contrary, that
\(s=\beta+it\ne1\) is a zero of \(\zeta\) with
\[
  R=|t|+2,\qquad L=\log R,\qquad d=\frac{2^{-20}}{L},
  \qquad \beta\ge 1-d .
\]
All logarithms of positive real numbers are natural logarithms.
Since \(L\ge\log 2>1/2\), we have \(0<d\le 1/100\).
The nonvanishing of \(\zeta\) on \(\Re z\ge 1\), away from its pole,
gives \(\beta<1\). In particular, \(s\ne0\) and \(F(s)=0\).

Since \(\beta\ge1-d\ge 99/100\), the height separation
in \eqref{eq:height-separation} applies. More precisely, its estimate at
this zero reads
\[
    \beta\le |s-1|\le 1-\beta+|t|,
    \qquad |t|\ge 2\beta-1\ge 1-2d\ge \tfrac12 .
\]
The last, weaker lower bound suffices to control both nonreal pole terms.

\subsection{One disk estimate at three centers}

Put \(a=\beta+5d\). Then \(4d\le a-1\), so \(1<a\le 1+5d\).
For real \(v\) with \(|v|\le 2|t|\), define
\[
 z_v=a+iv,\qquad f_v(w)=F(z_v+w),\qquad
 B_v=\frac{(16R)^2}{|F(z_v)|},\qquad
 H=1+8\log 2+2L .
\]
The denominator is positive because \(a>1\).
On \(|w|\le 1/2\) we have \(\Re(z_v+w)>1/2\). Also,
\[
   |z_v+w|\le a+|v|+\tfrac12\le \tfrac{31}{20}+2|t|,
   \qquad 2|z_v+w|+2\le \tfrac{51}{10}+4|t|\le 16R .
\]
Thus \lean{ClassicalZFR.growth} gives
\(|f_v(w)|\le(16R)^2=B_v|f_v(0)|\).
If \(f_v(w)=0\), then \(\Re(z_v+w)\le1\), so \(\Re w\le1-a<0\).
These are the growth and zero-location hypotheses of \eqref{eq:landau-disk},
proved by \lean{ClassicalZFR.landau}.

Define two real quantities by
\[
   D_v=\Re\sum_{n\ge 1}\frac{\Lambda(n)}{n^{z_v}},
   \qquad P_v=\Re\frac{1}{z_v-1},
\]
where \(\Lambda\) denotes the von Mangoldt function. For \(a>1\) the
Dirichlet series converges absolutely and equals \(-\zeta'(z_v)/\zeta(z_v)\).
Logarithmic differentiation of \(F\) consequently gives
\(\Re(f_v'(0)/f_v(0))=P_v-D_v\).
This is the sign convention in \lean{ClassicalZFR.good}. With \(x=0\),
or with an admissible zero witness of weight \(x\), the disk estimate reads
\begin{equation}
   x+D_v\le 24\bigl(H-\log|F(z_v)|\bigr)+P_v .
   \label{eq:three-window}
\end{equation}
Indeed, \(1+\log B_v=H-\log|F(z_v)|\). The term \(8\log 2\) is
\(\log(16^2)\), the coefficient of \(L\) is \(2\), and \(24\) multiplies
the whole expression.

Use \eqref{eq:three-window} at \(v=0,t,2t\). The outer two applications
take \(x=0\). At \(v=t\) the translated zero is \(\rho=-5d\), since
\(z_t-5d=s\). It lies in the witness disk because \(5d\le1/20<1/6\), and
its weight is \(-\Re(\rho^{-1})=1/(5d)\). The three pole terms satisfy
\begin{align}
 P_0&=\frac{1}{a-1}\le\frac{1}{4d},&
 P_t&\le\frac{1}{|z_t-1|}\le\frac{1}{|t|},&
 P_{2t}&\le\frac{1}{|z_{2t}-1|}\le\frac{1}{2|t|}.
 \label{eq:three-poles}
\end{align}
For the last two estimates we used \(\Re u\le |u|\) and
\(|z_v-1|\ge |v|\). In particular \(4P_t+P_{2t}\le9\), since
\(|t|\ge1/2\).

\subsection{The two weighted positivity estimates}

The first estimate controls the three center values inside the logarithms
in \eqref{eq:three-window}. The Dirichlet character lemma
\lean{norm_LFunction_product_ge_one} applies to the trivial character
of modulus \(1\), real displacement \(a-1>0\), and height \(t\). It gives
\[
       1\le |\zeta(z_0)|^3|\zeta(z_t)|^4|\zeta(z_{2t})|.
\]
At each center \(|z_v-1|\ge a-1\ge4d\).
Substituting \(|\zeta(z_v)|=|F(z_v)|/|z_v-1|\) and taking real
logarithms therefore gives
\begin{equation}
 \Sigma_F:=3\log|F(z_0)|+4\log|F(z_t)|+\log|F(z_{2t})|
             \ge 8\log(4d).
 \label{eq:three-product}
\end{equation}
The exponent \(8=3+4+1\) records all the factors introduced by clearing
the pole of \(\zeta\).

The second estimate concerns \(D_v\), the Dirichlet series in the
logarithmic derivative. Absolute convergence at real part \(a>1\) allows
us to take real parts term by term and combine the three sums. For
\(\theta_n=t\log n\) the result is
\begin{align}
 3D_0+4D_t+D_{2t}
   &=\sum_{n\ge1}\Lambda(n)n^{-a}
           \bigl(3+4\cos\theta_n+\cos(2\theta_n)\bigr)\notag\\
   &=\sum_{n\ge1}2\Lambda(n)n^{-a}(1+\cos\theta_n)^2\ge0 .
 \label{eq:three-series}
\end{align}
The double-angle identity gives the square, and every coefficient
\(\Lambda(n)n^{-a}\) is nonnegative. In the formal proof
\lean{HasSum.nonneg} sums these pointwise inequalities.
Thus \eqref{eq:three-product} gives a lower bound for a sum of logarithms,
while \eqref{eq:three-series} gives the sign of a sum of derivatives.

\subsection{Closing the numerical inequality}

Multiply \eqref{eq:three-window} at \(0,t,2t\) by \(3,4,1\),
respectively, and add. By \eqref{eq:three-series} the left side is at
least \(4/(5d)\). On the right, \eqref{eq:three-poles} bounds the total
pole contribution by \(3/(4d)+9\). Hence
\[
 \frac{4}{5d}
 \le 192H-24\Sigma_F+\frac{3}{4d}+9 .
\]
Inserting \eqref{eq:three-product} and using
\(4/5-3/4=1/20\) gives
\begin{equation}
     \frac{1}{20d}\le 192\bigl(H-\log(4d)\bigr)+9 .
     \label{eq:three-final}
\end{equation}
Here \(4d=2^{-18}/L\). In particular,
\(H-\log(4d)=1+26\log 2+2L+\log L\), with \(26=8+18\).
Substituting into \eqref{eq:three-final} yields
\[
    \frac{2^{20}}{20}L
        \le 201+4992\log 2+384L+192\log L .
\]
To check the constants, use \(\log L\le L\) and \(\log 2<1\).
The last inequality would imply
\[
     (2^{20}-11520)L
       \le 4020+99840\log 2 < 103860 .
\]
But \(2^{20}-11520=1037056\) and \(L\ge\log 2>1/2\), so its left
side exceeds \(518528\). This contradiction proves
\(\zeta(s)\ne0\) throughout the claimed region for \(s\ne1\), with the
explicit constant \(2^{-20}\).
